% ============================================================
% FORMAL FRAMEWORK — authoritative source for the theory section.
% For the Code session. This file consolidates the model developed in
% chat and SUPERSEDES the earlier Revised_Framework doc: where the two
% differ, this file governs. It is the model (§3, "A Constraint with
% Exit"); the institution's definitions (Scope/Height/Exit, R_j, L, the
% R x L match, the franchise cases) are in WhoMustAgree_section2_draft.tex
% (§2) and are assumed here. Place the two sections adjacently.
%
% BINDING CONSTRAINTS (do not violate; these are the data-audit findings):
%   1. Operationalise the legal menu L as government TYPE (ordinal, de jure,
%      predetermined). No graded de jure L^law exists in the data; the
%      continuous R_j x L^law match is a STATED EXTENSION, never a result.
%   2. op_referendum_strict is theta (the requirement), NOT L. The
%      R_j x theta comparison is the SUBSTITUTION test (P1/P3), not the exit
%      match. Never label it as the match; it is cross-state and descriptive.
%   3. P2 is a STRATIFIED RD reported as a schedule {tau_k} with honest CIs:
%      heterogeneity across non-random groups, NOT an identified causal effect
%      of exit. Do not give any single interaction coefficient a causal sign.
%   4. P3's aggregate claim needs an EQUIVALENCE test, not an insignificant
%      coefficient. Distinguish the extensive margin (whether a government
%      borrows at all: a tight null) from the aggregate amount (imprecise,
%      not to be read as zero).
%   5. R_j and the strata are PREDETERMINED, coded from function and statute,
%      never from realised avoidance behaviour.
% ============================================================

\section{Framework: A Constraint with Exit}

The institution of the previous section gives a government a choice rather than
a command. It must finance its investment somehow, and the consent requirement
governs only one of the routes available to it. This section formalises that
choice and states what it implies.

\subsection{The choice}

A government deciding to finance investment $j$ chooses a channel
$m \in \{V, E\}$: the voted channel $V$, on which authorisation must be sought,
or an exempt channel $E$, on which it need not. Conditional on undertaking the
project it selects the cheaper channel,
\[
m^{*} = \arg\min_{m \in \{V,E\}} C_m,
\]
and undertakes the project when its benefit covers that cost,
$B_j \geq \min\{C_V, C_E\}$.

The voted channel carries $C_V = C_V(\theta_s, e_i)$: the financing cost
together with the expected political and temporal cost of obtaining consent,
rising in the formal stringency $\theta_s$ and shaped by the electorate $e_i$
convened to grant it. The exempt channel carries $C_E = C_E(R_j, L_{is})$, low
where the project is chargeable ($R_j$ high) and a lawful exempt instrument
exists ($L_{is}$), and prohibitive where either is absent. The effective
outside option is the match between the two, $E_{ijs} = g(R_j, L_{is})$, with
economic feasibility and legal availability entering as complements.

\subsection{Two margins}

Raising the stringency of the requirement acts on two margins, whose relative
weight is set by the outside option. Where an exempt channel is cheap, a
stricter rule shifts financing onto it,
\[
\frac{\partial \Pr(E)}{\partial \theta} > 0,
\]
so that the requirement changes the instrument and the composition of what is
built rather than its quantity. Where exit is costly, the same increase acts on
quantity and timing: a project must obtain consent or wait. The response
therefore runs through instrument substitution, project substitution, and
quantity adjustment, in a mix governed by $E_{ijs}$.

\subsection{After a refusal}

A government that enters the voted channel and is refused chooses again, among
returning to the voters, rerouting through an exempt channel, or abandoning the
project. Rerouting is available where exit is cheap; where it is not, but
future consent remains feasible, refusal buys delay and resubmission;
abandonment is the residual, chosen only where both exit and renewed consent
are too costly to bear. A consent requirement therefore prices an adjustment
margin far more than it prevents the act.

\subsection{Incidence runs through two channels}

The requirement's incidence has two sources, and the analysis keeps them
separate. The principal channel is the outside option $E_{ijs}$: the
requirement binds where exit is poor and falls slack where it is rich. A second
and distinct channel is the electorate $e_i$ that $C_V$ convenes, whose
composition can make consent harder or easier to assemble; this channel is
treated as exploratory. The two must not be conflated, since they enter at
different points, $E$ through the exempt channel and $e_i$ through the cost of
the voted one. That the legal menu $L$ is itself politically constituted, and
in some forms restores a property-weighted franchise struck from the general
channel (Salyer, Ball v. James), is what makes the exit channel a political
institution rather than a financing detail.

\subsection{Predictions}

\paragraph{P1 (substitution).}
Where the requirement is more stringent, financing shifts away from the voted
instrument and toward exempt forms, and the composition of investment tilts
toward chargeable purposes, wherever a feasible legal exit exists. The test is
the association of rule stringency with the general-obligation share and with
project composition, the $R_j \times \theta$ comparison. This is a substitution
test and not the exit match: $\theta$ is the requirement, not the outside
option, and the comparison is cross-state and descriptive.

\paragraph{P2 (incidence through exit).}
The causal effect of obtaining authorisation is larger where the outside option
is poorer. The test stratifies the close-election design by predetermined exit
availability, leading with the school versus general-purpose fork and reporting
the effect as a schedule $\{\tau_k\}$ with honest confidence intervals. The
ordering across strata characterises heterogeneity across non-randomly assigned
groups; it does not identify the causal effect of exit itself, and no single
interaction coefficient is given a causal sign.

\paragraph{P3 (composition before quantity).}
Where substitution is available, stringency changes the financing and
composition of investment before it changes the aggregate. The composition
margins are estimated directly. The aggregate claim is supported by an
equivalence test rather than by an insignificant coefficient, and distinguishes
whether a government borrows at all, the extensive margin, from how much it
borrows, the aggregate amount, which is estimated with less precision.

% ============================================================
% GUIDANCE FOR ADJACENT SECTIONS (not part of §3; place as noted)
%
% EMPIRICAL STAGES (results architecture):
%   Stage I  — pre-vote response: rule -> voted share, GO share, composition,
%              and (extensive) total borrowing. Cross-state, descriptive
%              (this is R1/R2; R1 = institutional validation, R2 = the P1/P3
%              substitution test). NOT a first stage.
%   Stage II — close elections: the RD effect of authorisation on borrowing
%              and capital. Foreground the +11.5 out-of-design agreement here
%              as "two independent designs, one answer."
%   Stage III— where consent binds: the {tau_k} schedule, lead with the
%              predetermined type fork (P2).
%   Stage IV — response to refusal: delay, not denial.
%
% INFERENTIAL HIERARCHY (keep as its own short section):
%   RD-causal (incl. within predetermined strata) / descriptive association
%   (Stage I) / dynamic implication (persistence). State the boundaries before
%   they are attacked.
%
% CONTRIBUTION (three points only; restore the who-pays clause):
%   (1) constraints operate over choice sets: incidence cannot be read from
%       formal stringency alone;
%   (2) the outside option is an economics x institutions match (R_j x L),
%       so a formally neutral rule has a structured, predictable incidence;
%   (3) the design observes the behavioural response function, not only whether
%       the constrained act occurs.
%   Close on distribution: the requirement sorts who is governed AND who bears
%   the cost -- the ratepayers and users who service the exempt debt with no
%   vote. That who-pays clause is the distributive half and must be present.
% ============================================================
